\pdfbookmark{Общая характеристика работы}{characteristic}             % Закладка pdf
\section*{Общая характеристика работы}

\newcommand{\actuality}{\pdfbookmark[1]{Актуальность}{actuality}\underline{\textbf{\actualityTXT}}}
\newcommand{\progress}{\pdfbookmark[1]{Разработанность темы}{progress}\underline{\textbf{\progressTXT}}}
\newcommand{\aim}{\pdfbookmark[1]{Цели}{aim}\underline{{\textbf\aimTXT}}}
\newcommand{\tasks}{\pdfbookmark[1]{Задачи}{tasks}\underline{\textbf{\tasksTXT}}}
\newcommand{\aimtasks}{\pdfbookmark[1]{Цели и задачи}{aimtasks}\aimtasksTXT}
\newcommand{\novelty}{\pdfbookmark[1]{Научная новизна}{novelty}\underline{\textbf{\noveltyTXT}}}
\newcommand{\influence}{\pdfbookmark[1]{Практическая значимость}{influence}\underline{\textbf{\influenceTXT}}}
\newcommand{\methods}{\pdfbookmark[1]{Методология и методы исследования}{methods}\underline{\textbf{\methodsTXT}}}
\newcommand{\defpositions}{\pdfbookmark[1]{Положения, выносимые на защиту}{defpositions}\underline{\textbf{\defpositionsTXT}}}
\newcommand{\reliability}{\pdfbookmark[1]{Достоверность}{reliability}\underline{\textbf{\reliabilityTXT}}}
\newcommand{\probation}{\pdfbookmark[1]{Апробация}{probation}\underline{\textbf{\probationTXT}}}
\newcommand{\sootvetstvie}{\pdfbookmark[1]{Соответствие паспорту специальности}{contribution}\underline{\textbf{\sootvetstvieTXT}}}
\newcommand{\contribution}{\pdfbookmark[1]{Личный вклад}{contribution}\underline{\textbf{\contributionTXT}}}
\newcommand{\publications}{\pdfbookmark[1]{Публикации}{publications}\underline{\textbf{\publicationsTXT}}}

\input{common/characteristic} % Характеристика работы по структуре во введении и в автореферате не отличается (ГОСТ Р 7.0.11, пункты 5.3.1 и 9.2.1), потому её загружаем из одного и того же внешнего файла, предварительно задав форму выделения некоторым параметрам

%Диссертационная работа была выполнена при поддержке грантов \dots

%\underline{\textbf{Объем и структура работы.}} Диссертация состоит из~введения,
%четырех глав, заключения и~приложения. Полный объем диссертации
%\textbf{ХХХ}~страниц текста с~\textbf{ХХ}~рисунками и~5~таблицами. Список
%литературы содержит \textbf{ХХX}~наименование.

\pdfbookmark{Содержание работы}{description}                          % Закладка pdf
\section*{Содержание работы}

В диссертационной работе решена актуальная научно-техническая задача повышения энергетической эффективности осевых вентиляторов за счёт применения меридионального ускорения потока и разработки усовершенствованного метода их проектирования.
Разработаны математические и инженерные подходы, позволяющие учитывать влияние геометрии проточной части на распределение параметров потока, а также проведены численные и экспериментальные исследования, подтвердившие достоверность предложенных решений.

Во \underline{\textbf{введении}} обосновывается актуальность
исследований, проводимых в~рамках данной диссертационной работы,
приводится обзор научной литературы по~изучаемой проблеме,
формулируется цель, ставятся задачи работы, излагается научная новизна
и практическая значимость представляемой работы.

\underline{\textbf{Первая глава}} посвящена анализу современных методов проектирования осевых вентиляторов, а также существующих подходов к моделированию течения в решётках профилей.
Установлено, что традиционные методики проектирования, основанные на цилиндрической форме проточной части, не учитывают в полной мере влияние меридионального ускорения потока и тем самым ограничивают возможности повышения напора и КПД вентиляторов.

Существующие методы учёта влияния формы проточной части основываются на решении дифференциальных уравнений течения сеточными методами на электронно-вычислительной технике.
Такой подход к моделированию трудоёмок и требователен к вычислительной мощности, что делает его не актуальным при проведении оптимизации проточной части и профилировании лопаток.
Этим обосновывается целесообразность разработки нового метода, позволяющего описывать течение на осесимметричных поверхностях тока переменной толщины и учитывать изменение формы проточной части в меридиональном сечении.

Во \underline{\textbf{второй главе}} разработана математическая модель и реализован метод моделирования течения невязкой жидкости на осесимметричных поверхностях тока переменной толщины, основанная на методе дискретных вихрей и конформном отображении поверхности тока.

Осесимметричные поверхности тока определяются по решению уравнения радиального равновесия в проточной части.
\[
	\frac{1}{\rho }\frac{\partial p}{\partial n}=\frac{c_{u}^{2}}{r}\cos\psi+\frac{c_{m}^{2}}{r_{m}}, \qquad
	\label{eq:radialEq}
\]
\begin{figure}[htb]
	\centering{
		\includegraphics[scale=0.75]{4763.docx.tmp/word/media/image5.png}}
	\caption{Линии тока и ортогонали в меридиональной плоскости проточной части в первом приближении.}
	\label{mer_setch}
\end{figure}

Конформное отображение используется для формирования расчётной области на комплексной плоскости.
\[
		x=\int \limits_{a_{0}}^{a}\frac{da}{r}\sqrt{1+{r_{a}^{\prime}}^{2}};
		y=-\varphi.
		\label{eq:conformalMapping}
\]

\begin{figure}[htb]
	\centering{
		\includegraphics[scale=0.7]{4763.docx.tmp/word/media/image2.png}}
	\caption{Конформное отображение осесимметричной поверхности на комплексную плоскость \(x0y\).}
	\label{fig:conformalMapping}
\end{figure}

На основе разработанного метода моделирования течения предложен быстрый метод профилирования лопаток для предварительного проектирования турбомашин.
Время определения геометрии контура лопатки в одном кольцевом сечении по высоте составляет порядка 1 с.
В отличие от существующих решений, профилирование осуществляется на произвольной осесимметричной поверхности и подходит для проточной части любого типа - осевой, диагональной и радиальной.

В \underline{\textbf{третьей главе}} выполнено численное (CFD) моделирование и экспериментальные исследования для проверки разработанного метода.

В численных расчётах RANS использованы модели турбулентности \(k\) ⁣−\(\omega\) проведено исследование влияния параметров сетки, граничных условий и модели турбулентности на результаты моделирования.

Выполнена верификация разработанного метода расчёта: расхождения между расчётными и экспериментальными характеристиками по коэффициенту напора на номинальном режиме не превышают 3\%, по коэффициенту полезного действия "--- 0.1\%.
\begin{figure}[htb]
	\centering
	\includegraphics[scale=0.85]{img/stand}
	\caption{Схема стенда.}
	\label{stend}
\end{figure}

\begin{figure}[htb]
	\centering
	\includegraphics[scale=0.5]{img/characteristic}
	\caption{Аэродинамическая характеристика вентилятора.}
	\label{charact}
\end{figure}

Получено удовлетворительное совпадение распределений скоростей и давлений, что подтверждает достоверность математической модели и корректность применённой методики.
\begin{figure}[htb]
	\centering
	\includegraphics[scale=0.5]{img/c2m}
	\caption{Параметры потока по высоте проточной части.}
	\label{c2m}
\end{figure}

В \underline{\textbf{четвёртой главе}} математическим моделированием исследовано влияние формы проточной части и  меридионального ускорения потока на аэродинамические характеристики вентиляторов.

Выявлены особенности течения, такие как влияние формы поверхности тока, изменение Hт при сужении проточной части и при вращении решётки. Предложены способы оценки их влияния;

Показано влияние меридионального ускорения на возможность повышения Hт для элементарного рабочего колеса, элементарной и полной ступени вентилятора с меридиональным ускорением потока.

\underline{\textbf{В пятой главе}} приведены примеры применения разработанного метода при проектировании вентиляторов новой схемы.

Проведено сравнение метода моделирования течения с CFD RANS.

Показана большая точность профилирования, по сравнению с существующими методами основанными на продувках плоских решёток профилей.

Проведено сравнение характеристик вентиляторов схемы РК+СА с меридиональным ускорением потока и без, полученных численным моделированием CFD RANS.

Приведены примеры профилирования лопаток различных устройств:
\begin{itemize}
\item Направляющий аппарат в S-образном канале
\item Рабочее колесо радиального вентилятора
\item Рабочее колесо центробежного компрессора
\end{itemize}




\FloatBarrier
\pdfbookmark{Заключение}{conclusion}                                  % Закладка pdf
В \underline{\textbf{заключении}} приведены основные результаты работы, которые заключаются в следующем:
%% Согласно ГОСТ Р 7.0.11-2011:
%% 5.3.3 В заключении диссертации излагают итоги выполненного исследования, рекомендации, перспективы дальнейшей разработки темы.
%% 9.2.3 В заключении автореферата диссертации излагают итоги данного исследования, рекомендации и перспективы дальнейшей разработки темы.


\textit{В заключении подводятся итоги работы.
Формулируются основные выводы (7-9) по результатам исследований.
Приводятся сведения об апробации, полноте опубликования в научной печати основного содержания диссертации, ее результатов, выводов.
Приводятся сведения о защищенности технических решений авторскими свидетельствами (патентами).
Указываются предприятия, где внедрены результаты диссертационной работы и где еще они могут быть использованы.
Этот раздел занимает до восьми страниц текста.
Можно построить заключение к диссертации по схеме выполнения общей характеристики работы, приводимой в автореферате, что позволит усилить единство диссертации и автореферата.}

В диссертационной работе решена актуальная научно-техническая задача повышения коэффициента напора осевых вентиляторов при сохранении их энергетической эффективности за счёт применения меридионального ускорения потока и разработки методов их проектирования.

\begin{enumerate}
  \item Разработана и реализована математическая модель квазитрёхмерного течения в лопаточных венцах вентиляторов с меридиональным ускорением, основанная на модификации метода дискретных вихрей и конформном отображении осесимметричных поверхностей тока переменной толщины.
  Результатом численного моделирования является поле скоростей и давления в межлопаточном канале вращающихся и неподвижных пространственных решёток профилей.

  Основное отличие от существующих решений заключается в том, что модель основана на особенностях потенциального течения, что определяет её преимущества:
 \begin{enumerate}
  \item не требуется построения расчётной сетки и задания начального приближения;

  \item моделирование сводится к решению СЛАУ с относительно небольшим количеством неизвестных (примерно 150);

  \item результаты моделирования при различных режимах течения являются линейнозависимыми, течение на любом режиме работы решётки профилей моделируется как линейная комбинация результатов двух известных режимов.
\end{enumerate}

  \item На базе модели течения разработан метод профилирования лопаток вентилятора с меридиональным ускорением потока, осевых, диагональных и центробежных турбомашин.
  Использование теоретических характеристик решёток профилей позволило распространить подходы к профилированию осевых турбомашин на машины с произвольными меридиональными формами.
  Проведённые проверочные RANS-расчёты продемонстрировали удовлетворительное совпадение распределения давления на контурах лопаток, а так же расчётных коэффициентов теоретического напора в пределах 3,5\% до уточнения влияния вязкости, и менее 1\% после уточнения.

  \item Проведённое экспериментальное исследование характеристик модели вентилятора с меридиональным ускорением, спрофилированного разработанным методом, позволило подтвердить расчётные интегральные характеристики, такие как коэффициент теоретического напора \(\bar{H}_\text{т}\), и распределения \(\bar{H}_\text{т}\) и меридиональных компонент скорости \(\bar{c}_2m\) по высоте проточной части.
  Разница между расчётным и экспериментальным значениями \(\bar{H}_\text{т}\) составила менее 1\% на номинальном режиме и менее 3\% на остальной исследованной части рабочей характеристики.

  \item С использованием разработанных методов математического моделирования течения, было исследование влияние формы проточной части на характеристики вентиляторов с меридиональным ускорением, подтверждённые RANS-расчётами и экспериментом:
  \begin{enumerate}
  \item Показана зависимость коэффициента теоретического напора \(\bar{H}_\text{т}\) от сужения проточной части и  предложен способ оценки его влияния.
  В сравнении с экспериментом ошибка оценки \(\bar{H}_\text{т}\) составила менее 5\% при расширении проточной части и менее 3\% при сужении;

  \item Показано влияние меридионального ускорения на возможность повышения коэффициента теоретического напора \(\bar{H}_\text{т}\) для элементарного рабочего колеса, элементарной и полной ступени.
  Показано, что при сохранении КПД ступени, \(\bar{H}_\text{т}\) может быть значительно увеличен.
  В зависимости от относительного диаметра втулки и отношения динамического напора ступени к полному, \(\bar{H}_\text{т}\) может быть увеличен более чем в два раза, в сравнении с цилиндрической ступенью.
  Это позволяет при фиксированных значениях напора, расхода и мощности снизить частоту вращения ротора на 30\%, что приводит к снижению уровня суммарной звуковой мощности на 4,65~дБ (в 2,92 раз).
  \end{enumerate}

\end{enumerate}


\pdfbookmark{Литература}{bibliography}                                % Закладка pdf
%При использовании пакета \verb!biblatex! список публикаций автора по теме диссертации формируется в разделе <<\publications>>\ файла \verb!common/characteristic.tex!  при помощи команды \verb!\nocite!

\ifdefmacro{\microtypesetup}{\microtypesetup{protrusion=false}}{} % не рекомендуется применять пакет микротипографики к автоматически генерируемому списку литературы
\urlstyle{rm}                               % ссылки URL обычным шрифтом
\ifnumequal{\value{bibliosel}}{0}{% Встроенная реализация с загрузкой файла через движок bibtex8
    \renewcommand{\bibname}{\large \bibtitleauthor}
    \nocite{*}
    \insertbiblioauthor           % Подключаем Bib-базы
    %\insertbiblioexternal   % !!! bibtex не умеет работать с несколькими библиографиями !!!
}{% Реализация пакетом biblatex через движок biber
    % Цитирования.
    %  * Порядок перечисления определяет порядок в библиографии (только внутри подраздела, если `\insertbiblioauthorgrouped`).
    %  * Если не соблюдать порядок "как для \printbibliography", нумерация в `\insertbiblioauthor` будет кривой.
    %  * Если цитировать каждый источник отдельной командой --- найти некоторые ошибки будет проще.
    %
    %% authorvak
    \nocite{vakbib1}%
    \nocite{vakbib2}%
    \nocite{furaschov2022}%
    \nocite{furaschov2022a}%
    \nocite{furashov2025}%
    \nocite{furashov2026}%
    %
    %% authorwos
    \nocite{wosbib1}%
    %
    %% authorscopus
    \nocite{scbib1}%
    %
    %% authorpathent
    %\nocite{patbib1}%
    %
    %% authorprogram
    %\nocite{progbib1}%
    %
    %% authorconf
    \nocite{confbib1}%
    \nocite{confbib2}%
    %
    %% authorother
    \nocite{bib1}%
    \nocite{bib2}%

    \ifnumgreater{\value{usefootcite}}{0}{
        \begin{refcontext}[labelprefix={}]
            \ifnum \value{bibgrouped}>0
                \insertbiblioauthorgrouped    % Вывод всех работ автора, сгруппированных по источникам
            \else
                \insertbiblioauthor      % Вывод всех работ автора
            \fi
        \end{refcontext}
    }{
        \ifnum \totvalue{citeexternal}>0
            \begin{refcontext}[labelprefix=A]
                \ifnum \value{bibgrouped}>0
                    \insertbiblioauthorgrouped    % Вывод всех работ автора, сгруппированных по источникам
                \else
                    \insertbiblioauthor      % Вывод всех работ автора
                \fi
            \end{refcontext}
        \else
            \ifnum \value{bibgrouped}>0
                \insertbiblioauthorgrouped    % Вывод всех работ автора, сгруппированных по источникам
            \else
                \insertbiblioauthor      % Вывод всех работ автора
            \fi
        \fi
        %  \insertbiblioauthorimportant  % Вывод наиболее значимых работ автора (определяется в файле characteristic во второй section)
        \begin{refcontext}[labelprefix={}]
            \insertbiblioexternal            % Вывод списка литературы, на которую ссылались в тексте автореферата
        \end{refcontext}
        % Невидимый библиографический список для подсчёта количества внешних публикаций
        % Используется, чтобы убрать приставку "А" у работ автора, если в автореферате нет
        % цитирований внешних источников.
        \printbibliography[heading=nobibheading, section=0, env=countexternal, keyword=biblioexternal, resetnumbers=true]%
    }
}
\ifdefmacro{\microtypesetup}{\microtypesetup{protrusion=true}}{}
\urlstyle{tt}                               % возвращаем установки шрифта ссылок URL
